
\section{Proof of Theorem \ref{thm:master_W2}}\label{sec proof thm 1}

\noindent\textbf{Notation.}~For the reader's convenience, we recall some necessary notation from Section \ref{sec main}.  We say that a function $f:\R^m\rightarrow\R$ is pseudo-Lipschitz of order $k$, denoted $f\in\rm{PL}(k)$, if there is a constant $L>0$ such that for all $\x,\y\in\R^m$,
$
|f(\x) - f(\y)|\leq L(1+\tn{\x}^{k-1}+\tn{\y}^{k-1})\|\x-\y\|_2 
$
(See also Section \ref{SM useful fact}). 
We say that a sequence of probability distributions $\nu_p$ on $\R^m$ {converges in $W_k$} to 
$\nu$, written $\nu_p\stackrel{W_k}{\Longrightarrow} \nu$, if $W_k(\nu_p,\nu) \rightarrow 0$ as $p \rightarrow \infty$. An equivalent definition is that, for any $f\in\rm{PL}(k)$, $\lim_{p\rightarrow}\E f(X_p)=\E f(X)$, where expectation is with respect to $X_p\sim\nu_p$ and $X\sim\nu$ (e.g., \cite{montanari2017estimation}).
%
For a sequence of random variables $\mathcal{X}_{p}$ that converge in probability to some constant $c$ in the limit of Assumption \ref{ass:linear} below, we  write $\mathcal{X}_{p}\rP c$. For a sequence of event $\Ec_p$ for which $\lim_{p\rightarrow}\Pro(\Ec_p) = 1$, we say that $\Ec_p$ occurs \emph{with probability approaching 1}. For this, we will often use the shorthand ``wpa. 1".


\vspace{5pt}
Let $\X\in\R^{n\times p}$ have zero-mean and normally distributed rows with a diagonal covariance matrix $\bSi=\E[\x\x^T]$. Given a ground-truth vector $\betas$ and labels $\y=\X\betas+\sigma \z,~\z\sim\Nn(0,\Iden_n)$, we consider the least-squares problem subject to the minimum Euclidian norm constraint (as $\kappa=p/n>1$) given by
\begin{align}\label{eq:PO_beta}
\min_{\bt}\frac{1}{2}\tn{\bt}^2\quad\text{subject to}\quad \y=\X\bt.
\end{align}
%, we solve the least-squares problem  \eqref{optim me2} subject to the . 
It is more convenient to work with the following change of variable: $\w:=\sqrt{\Sigmab}(\bt-\betas)$. With this, the optimization problem  in \eqref{eq:min_norm} can be rewritten as
\begin{align}\label{eq:PO}
\Phi(\X)=\min_{\w} \frac{1}{2}\tn{\Sigmab^{-1/2}\w+\betas}^2\quad\text{subject to}\quad \Xb\w=\sigma \z,
\end{align}
where we write $\Xb=\X\Sigmab^{-1/2}\distas\Nn(0,1)$. First, using standard arguments, we show that the solution of \eqref{eq:PO} is bounded. Hence, we can constraint the optimization in a sufficiently large compact set without loss of generality.

\begin{lemma}[Boundedness of solution]\label{lem:bd_PO}
Let $\hat\w_n:=\hat\w_n(\X,\z)$ be the minimizer in \eqref{eq:PO}. Then, with probability approaching 1, it holds that $\hat\w_n\in\Bc$, where
$$\Bc:=\left\{\w\,|\,\|\w\|_2\leq B_{+} \right\},\qquad B_+:=4\sqrt{\Sigma_{\max}}\frac{2(\sqrt{\kappa}+1)\sqrt{\Sigma_{\max}} \sqrt{\E\left[B^2\right]} + \sigma}{\sqrt{\Sigma_{\min}}(\sqrt{\kappa}-1)} + \sqrt{\Sigma_{\max}}\sqrt{\E\left[B^2\right]}.
$$ 
\end{lemma}
\begin{proof}
First, we show that the min-norm solution $\hat\betab=\X^T(\X\X^T)^{-1}\y$ of \eqref{eq:PO_beta} is bounded. Here, we used the fact that $\kappa>1$, thus $\X\X^T$ is invertible wpa 1. We
have,
\begin{align}
\tn{\hat\betab_n}^2 = \y^T(\X\X^T)^{-1}\y \leq \frac{\tn{\y}^2}{\la_{\min}(\X\X^T)}  = \frac{\tn{\y}^2}{\la_{\min}(\Xb\Sigmab\Xb^T)} \leq  \frac{\tn{\y}^2}{\la_{\min}(\Xb\Xb^T)\,\Sigma_{\min}} = \frac{\tn{\y}^2}{\sigma_{\min}^2(\Xb)\,\Sigma_{\min}}. \label{eq:Ubb}
\end{align}
But, wpa 1, 
$
\sigma_{\min}(\Xb)/\sqrt{n} \geq \frac{1}{2}\left(\sqrt{\kappa}-1\right).
$
Furthermore, 
$
\|\y\|_2 \leq \|\Xb\Sigmab^{1/2}\betas\|_2 + \sigma\|\z\|_2 \leq \sigma_{\max}(\Xb)\sqrt{\Sigma_{\max}} \|\betas\|_2 + \sigma\|\z\|_2.
$
Hence, wpa 1,
$$
\|\y\|_2/\sqrt{n} \leq 4(\sqrt{\kappa}+1)\sqrt{\Sigma_{\max}} \sqrt{\E\left[B^2\right]} + 2\sigma,
$$
where we used the facts that wpa 1: $\|z\|_2/\sqrt{n}\rP 1$, $\sigma_{\max}(\Xb)<2(\sqrt{\kappa}+1)$ and by Assumption \ref{ass:mu}:
$$
\|\betas\|_2^2 = \frac{1}{p}\sum_{i=1}^{p}(\sqrt{p}\betas_i)^2 \rP \E\left[B^2\right].
$$
Put together in \eqref{eq:Ubb}, shows that
\begin{align}
\tn{\hat\betab_n} < \frac{4(\sqrt{\kappa}+1)\sqrt{\Sigma_{\max}} \sqrt{\E\left[B^2\right]} + 2\sigma}{\sqrt{\Sigma_{\min}}(\sqrt{\kappa}-1)/2} =: \tilde{B}_+.\label{eq:bd_beta}
\end{align}
Recalling that $\hat\w_n= \sqrt{\Sigmab}\betab - \sqrt{\Sigmab}\betas$, we conclude, as desired, that wpa 1,
$
\tn{\hat\w_n} \leq \sqrt{\Sigma_{\max}}\tilde{B}_+ + \sqrt{\Sigma_{\max}} \sqrt{\E[B^2]} =: B_+.
$
%{\color{red}Missing proof}
\end{proof}

Lemma \ref{lem:bd_PO} implies that nothing changes in \eqref{eq:PO} if we further constrain $\w\in\Bc$ in \eqref{eq:PO}. Henceforth, with some abuse of notation, we let
\begin{align}\label{eq:PO_bd}
\Phi(\X):=\min_{\w\in\Bc} \frac{1}{2}\tn{\Sigmab^{-1/2}\w+\betas}^2\quad\text{subject to}\quad \Xb\w=\sigma \z,
\end{align}

Next, in order to analyze the primary optimization (PO) problem in \eqref{eq:PO_bd} in apply the CGMT \cite{thrampoulidis2015lasso}. Specifically, we use the constrained formulation of the CGMT, Theorem \ref{thm closed}.  Specifically, the auxiliary problem (AO) corresponding to \eqref{eq:PO_bd} takes the following form with $\g\sim\Nn(0,\Iden_n)$, $\h\sim\Nn(0,\Iden_p)$, $h\sim \Nn(0,1)$:% (non-rigorously)
\begin{align}
\phi(\g,\h) = \min_{\w} \frac{1}{2}\tn{\Sigmab^{-1/2}\w+\betas}^2\quad\text{subject to}\quad \tn{\g}\tn{\w~{\sigma}}\leq \h^T\w+\sigma h.\label{eq:AO_con}
\end{align}%\tn{\h}
%Set $\bar{\h}=\h/\sqrt{p}$.

We will prove the following result about the AO problem.
\begin{lemma}[Properties of the AO -- Overparameterized regime]\label{lem:AO}
Let $\phi_n=\phi(\g,\h)$ be the optimal cost of the minimization in \eqref{eq:AO_con}. Define $\bar\phi$ as the optimal cost of the following deterministic min-max problem
\begin{align}\label{eq:AO_det}
\bar\phi:=\max_{u\geq 0}\min_{\tau>0}~ \Dc(u,\tau):=\frac{1}{2}\left({u\tau} + \frac{u\sigma^2}{\tau} - u^2\kappa\,\E\left[ \frac{1}{\Lambda^{-1}+\frac{u}{\tau}} \right] - \E\left[\frac{N^2}{\Lambda^{-1}+\frac{u}{\tau}}\right] \right).
\end{align}
The following statements are true.

\noindent{(i).}~The AO minimization in \eqref{eq:AO_con} is $\frac{1}{\Sigma_{\max}}$-strongly convex and has a unique minimizer $\w_n:=\w_n(\g,\h)$.

\noindent{(ii).}~In the limit of $n,p\rightarrow\infty, p/n=\kappa$, it holds that $\phi(\g,\h)\rP\bar\phi$, i.e., for any $\eps>0$:
$$
\lim_{n\rightarrow\infty}\P\left(|\phi(\g,\h)-\bar\phi|>\eps\right) = 0.
$$

\noindent{(iii).} The max-min optimization in \eqref{eq:AO_det} has a unique saddle point $(u_*,\tau*)$ satisfying the following: 
$$
u_*/\tau_* = \xi\quad\text{and}\quad\tau_* = \gamma,
$$
where $\xi, \gamma>0$ are defined in Definition \ref{def:Xi}.

\noindent{(iv).}~Let $f:\R\rightarrow\R$ be a $\rm{PL}({k})$ function. Let $\betab_n=\Sigmab^{-1/2}\w_n + \betas$. Then,
$$
\frac{1}{p}\sum_{i=1}^{p}f\left(\betab_n(\g,\h),\betas,\Sigmab\right) \rP \E_{(B,\Lambda,H)\sim\mu\otimes	\Nn(0,1)}\left[f\left(X_{\kappa,\sigma^2}(B,\Lambda,H),B,\Lambda\right) \right].
$$

\end{lemma}
We prove Lemma \ref{lem:AO} in Section \ref{sec:proofAO}. Here, we show how this leads to the proof of Theorem \ref{thm:master_W2} when combined with the CGMT.



%\begin{lemma}\label{lem:deviation}
%Let $\hat\w_n=\hat\w_n(\X,\z)$ be a minimizer of the PO in \eqref{eq:PO}. 
%%Further let $\tau_n=\tau_n(\g,\h,\z)$ and $u_n=u_n(\g,\h,\z)$ be the unique saddle point of the min-max optimization in \eqref{eq:AO_4}. Define $\w_n=\w_n(\tau_n,u_n)$ as in \eqref{eq:w_n}. 
%For any pseudo-Lipschitz function $f:\R\rightarrow\R$ of order 2, it hold that
%\begin{align}
%\frac{1}{p} \sum_{i=1}^{p} f\left(\hat\w_{n,i}\right) \rP \E_\mu\left[X_{\kappa,\sigma^2}(\Lambda,N,H)  \right].
%\end{align}
%\end{lemma}


%\begin{proof}[Proof of Lemma \ref{lem:deviation}]


\noindent\textbf{Finishing the proof of Theorem \ref{thm:master_W2}:} For convenience, define 
$$F_n(\hat\betab,\betas,\Sigmab):=\frac{1}{p} \sum_{i=1}^{p} f\left(\sqrt{p}\hat\betab_{n,i},\sqrt{p}\betas_{i},\Sigmab_{ii}\right)\quad\text{and}\quad\alpha_*:=\E_\mu\left[f(X_{\kappa,\sigma^2}(\Lambda,B,H),B,\Lambda)\right].$$
Fix  any $\eps>0$. It suffices to prove that the solution $\hat\betab_n$ of the PO in \eqref{eq:PO_beta} satisfies $\hat\betab_n\not\in\Sc$ wpa. 1, where 
\begin{align}\label{eq:S_set}
\Sc = \Sc(\betas,\Sigmab) =\{\betab\bgl |F_n(\hat\betab_n,\betas,\Sigmab)-\alpha_*|\geq 2\eps\}.
\end{align}
%\ct{Note that this $\Sc$ is a random set now.}
In particular, define the ``perturbed" PO and AO problems (compare to \eqref{eq:PO} and \eqref{eq:AO_con}) as:
\begin{align}\label{eq:PO_S}
\Phi_S(\X)=\min_{\w\in\Sc} \frac{1}{2}\tn{\Sigmab^{-1/2}\w+\betas}^2\quad\text{subject to}\quad \Xb\w=\sigma \z,
\end{align}
and
\begin{align}
\phi_S(\g,\h)=\min_{\w\in\Sc} \frac{1}{2}\tn{\Sigmab^{-1/2}\w+\betas}^2\quad\text{subject to}\quad \tn{\g}\tn{\w~{\sigma}}\leq \h^T\w+\sigma h.\label{eq:AO_S}
\end{align}%\tn{\h}
where recall that $\Xb=\X\Sigmab^{-1/2}\distas\Nn(0,1)$, $\g\sim\Nn(0,\Iden_n)$, $\h\sim\Nn(0,\Iden_p)$, $h\sim \Nn(0,1)$ and we have used the change of variables $\w:=\sqrt{\Sigmab}(\betab-\betas)$ for convenience. 

 Using \cite[Theorem 6.1(iii)]{thrampoulidis2018precise} it suffices to find costants $\bar\phi, \bar\phi_S$ and $\eta>0$ such that the following hold:
\begin{enumerate}
\item $\bar\phi_S \geq \bar\phi + 3\eta$,
\item $\phi(\g,\h) \leq \bar\phi + \eta$, with probability approaching 1,
\item $\phi_S(\g,\h) \geq \bar\phi_S - \eta$, with probability approaching 1.
\end{enumerate}
In what follows, we  explicitly find $\bar\phi, \bar\phi_S,\eta$ such that the three conditions above hold.

\vspace{5pt}
\noindent\underline{Satisfying Condition 2}: Recall the deterministic min-max optimization in \eqref{eq:AO_det}. Choose $\bar\phi=\Dc(u_*,\tau_*)$ be the optimal cost of this optimization. From Lemma \ref{lem:AO}(ii), $\phi(\g,\h)\rP\bar\phi$. Thus, for any $\eta>0$, with probability approaching 1:
\begin{align}\label{eq:phi_lim}
\bar\phi + \eta \geq \phi(\g,\h) \geq \bar\phi - \eta.
\end{align}
Clearly then, Condition 2 above holds for any $\eta>0$. 

\vspace{5pt}
\noindent\underline{Satisfying Condition 3}: Next, we will show that the third condition holds for appropriate $\bar\phi$. 
%Let $\tau_n=\tau_n(\g,\h,\z)$ and $u_n=u_n(\g,\h,\z)$ be the unique saddle point of the min-max optimization in \eqref{eq:AO_4}. Define 
Let $\w_n=\w_n(\g,\h)$ be the unique minimizer of \eqref{eq:AO_con} as per Lemma \ref{lem:AO}(i), i.e., $\frac{1}{2}\tn{\Sigmab^{-1/2}\w_n+\bt^\st}^2 = \phi(\g,\h)$. Again from Lemma \ref{lem:AO}, the minimization in \eqref{eq:AO_con} is $1/\Sigma_{\max}$-strongly convex in $\w$. Here, $\Sigma_{\max}$ is the upper bound on the eigenvalues of $\Sigmab$ as per Assumption \ref{ass:inv}. Thus, for any $\tilde\eps>0$ and any feasible $\w$ the following holds (deterministically):
\begin{align}\label{eq:sc}
\frac{1}{2}\tn{\Sigmab^{-1/2}\w+\bt^\st}^2 \geq \phi(\g,\h) + \frac{\tilde\epsilon^2}{2C},~\text{provided that}~ \|\w-\w_n(\g,\h)\|_2 \geq \tilde\epsilon.
\end{align}
Now, we argue that $\betab\in\Sc$ implies that $\|\w-\w_n(\g,\h)\|_2\geq \tilde\eps$ wpa 1, for 
appropriate value of $\tilde\eps$ and $\w=\sqrt{\Sigmab}(\betab-\betas)$. Consider any $\betab\in\Sc$. First, by definition in \eqref{eq:S_set},
$$
|F_n(\betab,\betas,\Sigmab)-\alpha_*| \geq 2\eps.
$$
Second, by Lemma \ref{lem:AO}(iv), with probability approaching 1,
$$
|F(\betab_n(\g,\h),\betas,\Sigmab) - \alpha_*| \leq \epsilon.
$$
Third, we will show that wpa 1, there exists universal constant $C>0$ such that
\begin{align}
|F_n(\betab_n(\g,\h),\betas,\Sigmab) - F_n(\betab,\betas,\Sigmab)| \leq C {\|\betab_{n}(\g,\h) - \betab\|_2}\label{eq:dev2show}.
\end{align}
Before proving \eqref{eq:dev2show}, let us argue how combining the above three displays shows the desired. Indeed, wpa. 1,
\begin{align*}
2\eps &\leq |F_n(\betab,\betas,\Sigmab)-\alpha_*| \leq |F_n(\betab_n(\g,\h),\betas,\Sigmab) - F_n(\betab,\betas,\Sigmab)| + |F_n(\betab_n(\g,\h),\betas,\Sigmab) - \alpha_*| \\
&\leq \epsilon + C \,\|\betab-\betab_n\|_2. \\
&\qquad\Longrightarrow \|\betab-\betab_n\|_2 \geq {\eps}/{C}=:\hat\eps\\
&\qquad\Longrightarrow \|\w-\w_n\|_2 \geq \hat\eps\sqrt{\Sigma_{\min}}=:\tilde\eps.
\end{align*}
In the last line above, we recalled that $\betab=\Sigmab^{-1/2}\w+\betas$ and $\Sigmab_{i,i}\geq\Sigma_{\min},~i\in[p]$ by Assumption \ref{ass:inv}.

From this and \eqref{eq:sc}, we find that wpa 1,
$
\frac{1}{2}\tn{\Sigmab^{-1/2}\w+\betab}^2 \geq \phi(\g,\h) + \frac{\tilde\epsilon^2}{2C},~\text{for all}~ \w\in\Sc.
$
Thus, 
\begin{align}
\phi_S(\g,\h) \geq \phi(\g,\h) + \frac{\tilde\epsilon^2}{2\Sigma_{\max}}.\nn
\end{align}
%\end{proof}
When combined with \eqref{eq:phi_lim}, this shows that
\begin{align}
\phi_S(\g,\h) \geq \bar\phi + \frac{\tilde\epsilon^2}{2\Sigma_{\max}} - \eta.
\end{align}
Thus, choosing $\bar\phi_S = \bar\phi + \frac{\tilde\epsilon^2}{2\Sigma_{\max}}$ proves the Condition 3 above. 

To complete the proof, let us now show \eqref{eq:dev2show}. Henceforth, $C$ is used to denote a universal constant whose value can change from line to line.
% by the pseudo-Lipschitz property, boundedness of $\betab_n$ and $\betab$ (see Lemma \ref{lem:bd_PO}) the following holds. 
To simplify notation, let $\ab_i=(\sqrt{p}\betab_{n,i}(\g,\h),\sqrt{p}\betas_i,\Sigmab_{i,i})$ and $\bb_i=(\sqrt{p}\betab_i,\sqrt{p}\betas_i,\Sigmab_{i,i})$, for $i\in[p]$. Then, using $f\in\rm{PL}(k)$,
\begin{align}
|F_n(\betab_n(\g,\h),\betas,\Sigmab) - F_n(\betab,\betas,\Sigmab)| &\leq \frac{L}{p}\sum_{i=1}^p (1+\max\{\|\ab_i\|_2^{k-1},\|\bb_i\|_2^{k-1}\})\|\ab_i-\bb_i\|_2\nn \\
&= \frac{L}{p}\sum_{i=1}^p \left(1+\max\{\|\ab_i\|_2^{k-1},\|\bb_i\|_2^{k-1}\}\right)\cdot|\sqrt{p}\betab_{n,i}(\g,\h) - \sqrt{p}\betab_i|\nn\\
&\leq L\left(1+\max\{\frac{1}{p}\sum_{i=1}^p \|\ab_i\|_2^{2k-2},\frac{1}{p}\sum_{i=1}^p\|\bb_i\|_2^{2k-2}\} \right)^{1/2}{\|\betab_{n}(\g,\h) - \betab\|_2}\label{eq:LipC}\,.
\end{align}
The last inequality above follows by applying Cauchy-Schwartz.  
To reach \eqref{eq:dev2show} from the above, it suffices to show that
\begin{align}
\max\{\frac{1}{p}\sum_{i=1}^p \|\ab_i\|_2^{2k-2},\frac{1}{p}\sum_{i=1}^p\|\bb_i\|_2^{2k-2}\} \leq C,\label{eq:leqC}
\end{align}
for some universal constant $C$ (that may depend on $k$, but not on $n,p$).
% $\frac{1}{p}\sum_{i=1}^p \|\ab_i\|_2^{2k-2}<\infty$ and $\frac{1}{p}\sum_{i=1}^p \|\bb_i\|_2^{2k-2}<\infty$ as $p\rightarrow\infty$. But, 
To prove \eqref{eq:leqC}, note that 
 using $a^{\ell_1}b^{\ell_2}c^{\ell_3}\leq a^{\ell}+b^{\ell}+c^\ell,~\ell=\ell_1+\ell_2+\ell_3$, for some constant $C=C(k)>0$ it holds:
%\ct{Why is this (or something like this true)?}
%\som{arithmetic - geometric mean inequality?}
\begin{align}
\frac{1}{p}\sum_{i=1}^p \|\bb_i\|_2^{2k-2} {\leq} C\left(\frac{1}{p}\sum_{i=1}^p |\sqrt{p}\betab_{i}|^{2k-2} + \frac{1}{p}\sum_{i=1}^p |\sqrt{p}\betas_{i}|^{2k-2} + \frac{1}{p}\sum_{i=1}^p |\Sigmab_{i,i}|^{2k-2} \right)\label{eq:bdmom}\,.
\end{align}
\cts{The terms $\sum_{i=1}^p |\sqrt{p}\betas_{i}|^{2k-2} $ and $ \frac{1}{p}\sum_{i=1}^p |\Sigmab_{i,i}|^{2k-2}$ are bounded by assumption}\ct{Need to add assumption that $ \frac{1}{p}\sum_{i=1}^p |\Sigmab_{i,i}|^{2k-2}<\infty$ and $\frac{1}{p}\sum_{i=1}^p (\sqrt{p}|\betas_{i}|^{2k-2} $. OK by just assuming that $B$ is Definition \ref{def:Xi} is bounded.}. {\color{red}Furthermore we need to argue that, 
\begin{align}
\frac{1}{p}\sum_{i=1}^p |\sqrt{p}\betab_{i}|^{2k-2} <\infty
\end{align}
}
%\frac{1}{d}\sum_{i=1}^p |\betab_{i}|^{2k-2}\leq \|\betab\|_2^{2k-2}<C'(k)$, where the last inequality follows from feasibility of $\betab$, i.e., $\w=\sqrt{\Sigmab}(\betab-\betas)\in\Bc$; see Lemma \ref{lem:bd_PO} and \eqref{eq:bd_beta}. 
These together with \eqref{eq:bdmom} show that $\frac{1}{p}\sum_{i=1}^p \|\bb_i\|_2^{2k-2}<C$ wpa 1. Similarly, we can argue for $\frac{1}{p}\sum_{i=1}^p \|\ab_i\|_2^{2k-2}$, completing the proof of \eqref{eq:leqC}.



\vspace{5pt}
\noindent\underline{Satisfying Condition 1:} To prove Condition 1, we simply pick $\eta$ to satisfy the following
\begin{align}
\bar\phi_S > \bar\phi + 3 \eta ~\Leftarrow~ \bar\phi + \frac{\tilde\epsilon^2}{2\Sigma_{\max}} - \eta \geq \bar\phi + 3 \eta ~\Leftarrow~ \eta \leq \frac{\tilde\epsilon^2}{8\Sigma_{\max}}.
\end{align}


This completes the proof of Theorem \ref{thm:master_W2}.

\subsection{Proof of Lemma \ref{lem:AO}}\label{sec:proofAO}

~~~~
~~~~

\vp
\noindent\underline{Proof of (i):}~Strong convexity of the objective function in \eqref{eq:PO} is easily verified by the second derivative test. Note here that we use Assumption \ref{ass:inv} that $\lai\leq\Sigma_{\max},~i\in[p].$ Uniqueness of the solution follows directly from strong convexity. \ct{Strictly speaking we might need to also argue existence, i.e., feasibility of the AO. An indirect way is to show feasibility using the CGMT, but it  seems unnecessarily complicated?}



\vspace{5pt}
\noindent\underline{Proof of (ii):}~Using Lagrangian formulation, the solution $\w_n$ to \eqref{eq:AO_con} is the same as the solution to the following:
\begin{align}
\left(\w_n,u_n\right) :=\arg\min_{\w\in\Bc}\max_{u\geq 0} ~\frac{1}{2}\tn{\Sigmab^{-1/2}\w+\betas}^2 + u \left( \sqrt{\tn{\w}^2+\sigma^2} \tn{\gba} - \sqrt{\kappa}\,{\hba}^T\w + \frac{\sigma h}{\sqrt{n}} \right)\label{eq:AO_2}
\end{align}
where we have: (i) set $\gba := \g/\sqrt{n}$ and $\hba:= \h/\sqrt{p}$; (ii) recalled that $p/n=\kappa$; and, (iii) used $\left(\w_n,u_n\right)$ to denote the optimal solutions in \eqref{eq:AO_2}. The subscript $n$ emphasizes the dependence of $\left(\w_n,u_n\right)$ on the problem dimensions. Also note that (even though not explicit in the notation) $\left(\w_n,u_n\right)$ are random variables depending on the realizations of $\gba,\hba$ and $h$.

 Notice that the objective function above is convex in $\w$ and linear (thus, concave) in $u$. Thus, strong duality holds and we can flip the order of min-max. Moreover, in order to make the objective easy to optimize with respect to $\w$, we use the following variational expression for the square-root term $\sqrt{\tn{\w}^2+\sigma^2}$:
$$
\tn{\gba}\sqrt{\tn{\w}^2+\sigma^2} = \tn{\gba}\cdot\min_{\tau\in[\sigma,\sqrt{\sigma^2+B_+^2}]} \left\{ \frac{\tau}{2} + \frac{\tn{\w}^2+\sigma^2}{2\tau} \right\} = \min_{\tau\in[\sigma,\sqrt{\sigma^2+B_+^2}]} \left\{ \frac{\tau\tn{\gba}^2}{2} + \frac{\sigma^2}{2\tau} + \frac{\tn{\w}^2}{2\tau} \right\},
$$
where $B$ is defined in Lemma \ref{lem:bd_PO}. For convenience define the constraint set for the variable $\tau$ as $\Tc':=[\sigma,\sqrt{\sigma^2+B_+^2}]$. For reasons to be made clear later in the proof (see proof of statement (iii)), we consider the (possibly larger) set:
\[
\Tc:=[\sigma,\max\{\sqrt{\sigma^2+B_+^2},2\tau_*\}]\,
\]
where $\tau_*$ is as in the statement of the lemma.

The above lead to the following equivalent formulation of \eqref{eq:AO_2}:
\begin{align}
\left(\w_n,u_n,\tau_n\right) = \max_{u\geq 0}\min_{\tau\in\Tc} ~ \frac{u\tau\tn{\gba}^2}{2} + \frac{u\sigma^2}{2\tau} + \frac{u\sigma h}{\sqrt{n}} + \min_{\w} \left\{ \frac{1}{2}\tn{\Sigmab^{-1/2}\w+\betas}^2 + \frac{u}{2\tau}\tn{\w}^2 - u \sqrt{\kappa}\,{\hba}^T\w \right\} .\label{eq:AO_3}
\end{align}
\ct{To be fully rigorous, need to show here that the unconstrained min over $\w$ is the same as the constrained $\w\in\Bc$.}
The minimization over $\w$ is easy as it involves a strongly convex quadratic function. The optimal $\w':=\w'(\tau,u)$ (for fixed $(\tau,u)$) is given by
\begin{align}\label{eq:w'}
\w':=\w'(\tau,u) = -\left(\Sigmab^{-1}+\frac{u}{\tau}\Iden\right)^{-1}\left(\Sigmab^{-1/2}\betas-u\sqrt{\kappa}\hba\right), 
\end{align}
and \eqref{eq:AO_3} simplifies to
\begin{align}
\left(u_n,\tau_n\right)=\max_{u\geq 0}\min_{\tau\in\Tc} ~ \frac{u\tau\tn{\gba}^2}{2} + \frac{u\sigma^2}{2\tau} + \frac{u\sigma h}{\sqrt{n}} - \frac{1}{2} \left(\Sigmab^{-1/2}\betas-u\sqrt{\kappa}\hba\right)^T\left(\Sigmab^{-1}+\frac{u}{\tau}\Iden\right)^{-1} \left(\Sigmab^{-1/2}\betas-u\sqrt{\kappa}\hba\right)\,=:\Rc(u,\tau) .\label{eq:AO_4}
\end{align}
It can be checked by direct differentiation and the second-derivative test that the objective function in \eqref{eq:AO_4} is strictly convex in $\tau$ and strictly concave in $u$ in the domain $\{(u,\tau)\in\R_+\times\R_+\}$. Thus, the saddle point $(u_n,\tau_n)$ is unique. Specifically, this implies that the optimal $\w_n$ in \eqref{eq:AO_3} is given by (cf. \eqref{eq:w'})
\begin{align}\label{eq:w_n}
\w_n=\w'(\tau_n,u_n) = -\left(\Sigmab^{-1}+\frac{u_n}{\tau_n}\Iden\right)^{-1}\left(\Sigmab^{-1/2}\betas-u_n\sqrt{\kappa}\hba\right).
\end{align}

In what follows, we characterize the high-dimensional limit of the optimal pair $(u_n,\tau_n)$ in the limit $n,p\rightarrow\infty,~p/n\rightarrow\kappa$.
%Now that we have reached a min-max problem over only scalar variables, we are ready to study its convergence properties in the limit $n,p\rightarrow\infty,~p/n\rightarrow\kappa$. 
We start by analyzing the (point-wise) convergence of $\Rc(u,\tau)$. 
For the first three summands in \eqref{eq:AO_4}, we easily find that
$$
\left\{\frac{u\tau\tn{\gba}^2}{2} + \frac{u\sigma^2}{2\tau} + \frac{u\sigma h}{\sqrt{n}}\right\}~ \rP~\left\{
\frac{u\tau}{2} + \frac{u\sigma^2}{2\tau} \right\}.
$$
Next, we study the fourth summand. First, note that
\begin{align}
(u\sqrt{\kappa}\hba)^T\left(\Sigmab^{-1}+\frac{u}{\tau}\Iden\right)^{-1}(u\sqrt{\kappa}\hba) &= u^2\kappa\,\frac{1}{p}\h^T\left(\Sigmab^{-1}+\frac{u}{\tau}\Iden\right)^{-1}\h \nn\\
&= u^2\kappa\,\frac{1}{p}\sum_{i=1}^{p}\frac{\h_i^2}{\lai^{-1}+\frac{u}{\tau}} \nn\\
&\rP u^2\kappa\,\E\left[ \frac{1}{\Lambda^{-1}+\frac{u}{\tau}} \right].
\end{align}
% the quadratic form
%\begin{align}
%(u\sqrt{\kappa}\hba)^T\left(\Sigmab^{-1}+\frac{u}{\tau}\Iden\right)^{-1}(u\sqrt{\kappa}\hba) = u^2\kappa\,\frac{1}{p}\h^T\left(\Sigmab^{-1}+\frac{u}{\tau}\Iden\right)^{-1}\h
%\end{align}
%concentrates around its expectation \cite{?} $u^2\kappa\,\tr\left(\left(\Sigmab^{-1}+\frac{u}{\tau}\Iden\right)\right)^{-1}.$
In the last line, $\Lambda$ is a random variable as in Definition \ref{def:Xi}. \cts{Also, we used Assumption \ref{ass:mu} together with the facts that $\h$ is independent of $\Sigmab$ and  that the function $(x_1,x_2)\mapsto f(x_1,x_2)=x_1^2(x_2^{-1}+u/\tau)^{-1}$ is $\rm{PL}(2)$ assuming $x_2$ is bounded.} \ct{Question: Is it immediate that the empirical distribution of $(\beta,\Sigmab,\h)$ converges in $W_k$ to $\mu\otimes\Nn(0,1)$ given that $(\beta,\Sigmab)$ converges to $\mu$ and $\h$ is independent???}
%
 Second, we find that
\begin{align}
(\betas)^T\Sigmab^{-1/2}\left(\Sigmab^{-1}+\frac{u}{\tau}\Iden\right)^{-1}\Sigmab^{-1/2}\betas &= \frac{1}{p}(\sqrt{p}\betas)^T\left(\Iden+\frac{u}{\tau}\Sigmab\right)^{-1}(\sqrt{p}\betas) \nn\\
&= \frac{1}{p}\sum_{i=1}^{p}\frac{\left(\sqrt{p}\betas_i/\sqrt{\lai}\right)^2}{\lai^{-1}+\frac{u}{\tau}}\nn\\
&\rP \E\left[\frac{B^2\Lambda^{-1}}{\Lambda^{-1}+\frac{u}{\tau}}\right].
\end{align}
Here, $\Lambda,B$ are random variables as in Definition \ref{def:Xi} and \cts{we also used Assumption \ref{ass:mu} together with the fact that the function $(x_1,x_2)\mapsto f(x_1,x_2)=x_1^2x_2^{-1}(x_2^{-1}+u/\tau)^{-1}$ is $\rm{PL}(2)$ assuming $x_2$ is bounded.}
Third, \cts{by independence of $(\betas, \Sigmab)$ from $\h$}
\begin{align}
(u\sqrt{\kappa}\hba)^T\left(\Sigmab^{-1}+\frac{u}{\tau}\Iden\right)^{-1}\Sigmab^{-1/2}\betas = u\sqrt{\kappa} \cdot \frac{1}{p}\sum_{i=1}
^p \frac{\h_i\Sigmab_{i,i}^{-1/2}(\sqrt{p}\betas_i)}{\Sigmab_{i,i}^{-1}+\frac{u}{\tau}\Iden} ~\rP~ 0.
\end{align}
Putting these together, the objective $\Rc(u,\tau)$ in \eqref{eq:AO_4} converges point-wise in $u,\tau$ to
\begin{align}
\Rc(u,\tau)\rP\Dc(u,\tau) := \frac{1}{2}\left({u\tau} + \frac{u\sigma^2}{\tau} - u^2\kappa\,\E\left[ \frac{1}{\Lambda^{-1}+\frac{u}{\tau}} \right] - \E\left[\frac{B^2\Lambda^{-1}}{\Lambda^{-1}+\frac{u}{\tau}}\right] \right).\label{eq:conv_pt}
\end{align}
Note that $\Rc(u,\tau)$ (and thus, $\Dc(u,\tau)$) is convex in $\tau$ and concave in $u$. Thus, the convergence in \eqref{eq:conv_pt} is in fact uniform (e.g., \cite{AG1982}) and we can conclude that 
\begin{align}\label{eq:Dc0}
\phi(\g,\h) \rP \max_{u\geq 0}\min_{\tau\in\Tc}~ \Dc(u,\tau).
\end{align}
and
\begin{align}\label{eq:Dc}
(u_n,\tau_n) \rP (u_*,\tau_*):=\arg\max_{u\geq 0}\min_{\tau\in\Tc}~ \Dc(u,\tau).
\end{align}
In the proof of statement (iii) below, we show that the saddle point of \eqref{eq:Dc0} is $(u_*,\tau_*)$. In particular, $\tau_*$ is strictly in the interior of $\Tc$, which combined with convexity implies that
$$
 \max_{u\geq 0}\min_{\tau\in\Tc}~ \Dc(u,\tau) =  \max_{u\geq 0}\min_{\tau>0}~ \Dc(u,\tau) =: \bar\phi.
$$
This, together with the first display above proves the second statement of the lemma.

\vspace{5pt}
\noindent\underline{Proof of (iii):} 
Next, we compute the saddle point $(u_*,\tau_*)$ by studying the first-order optimality conditions of the srtictly concave-convex $\Dc(u,\tau)$. Specifically, we consider unconstrained minimization over $\tau$ and we will show that the minimum is achieved in the strict interior of $\Tc$. Direct differentiation gives
\begin{subequations}
\begin{align}
{\tau} + \frac{\sigma^2}{\tau} - 2u\kappa\E\left[ \frac{1}{\Lambda^{-1}+\frac{u}{\tau}} \right] + \frac{u^2}{\tau}\kappa \E\left[ \frac{1}{\left(\Lambda^{-1}+\frac{u}{\tau}\right)^2}\right] + \frac{1}{\tau}\E\left[\frac{B^2\Lambda^{-1}}{\left(\Lambda^{-1}+\frac{u}{\tau}\right)^2}\right] &= 0, \label{eq:fo1}\\
{u} - \frac{u\sigma^2}{\tau^2} - \frac{u^3}{\tau^2}\kappa \E\left[ \frac{1}{\left(\Lambda^{-1} + \frac{u}{\tau}\right)^2}\right] - \frac{u}{\tau^2} \E\left[\frac{B^2\Lambda^{-1}}{\left(\Lambda^{-1}+\frac{u}{\tau}\right)^2}\right] &= 0,\label{eq:fo2}
\end{align}
\end{subequations}
Multiplying \eqref{eq:fo2}  with $\frac{\tau}{u}$ and adding to \eqref{eq:fo1} results in the following equation
\begin{align}
\tau = u\kappa\E\left[ \frac{1}{\Lambda^{-1}+\frac{u}{\tau}} \right] ~\Leftrightarrow ~
\E\left[ \frac{1}{(\frac{u}{\tau}\Lambda)^{-1}+1} \right] = \frac{1}{\kappa} \label{eq:tauu}\,.
\end{align}
Thus, we have found that the ratio $\frac{u_*}{\tau_*}$ is the unique solution to the equation in \eqref{eq:tauu}. Note that this coincides with the Equation \eqref{eq:ksi} that defines the parameter $\xi$ in Definition \ref{def:Xi}.  The fact that \eqref{eq:tauu} has a unique solution for all $\kappa>1$ can be easily seen as $F(x)=\E\left[ \frac{1}{(x\Lambda)^{-1}+1} \right], x\in\R_+$ has range $(0,1)$ and is strictly increasing (by differentiation). 


Thus, we call $\xi=\frac{u_*}{\tau_*}$. Moreover, multiplying \label{eq:fo2} with $u$ leads to the following equation for $\tau_*$:
\begin{align}
u_*^2 = \sigma^2\xi^2 + u_*^2 \xi^2  \kappa \E\left[ \frac{1}{\left(\Lambda^{-1} +\xi\right)^2}\right] + \xi^2  \E\left[ \frac{B^2\Lambda^{-1}}{\left(\Lambda^{-1} + \xi\right)^2}\right] ~\Rightarrow~ \tau_*^2 = \frac{\sigma^2 + \E\left[ \frac{B^2\Lambda^{-1}}{\left(\Lambda^{-1} + \xi\right)^2}\right]}{1-\xi^2\kappa \E\left[ \frac{1}{\left(\Lambda^{-1} +\xi\right)^2}\right]} =  \frac{\sigma^2 + \E\left[ \frac{B^2\Lambda^{-1}}{\left(\Lambda^{-1} + \xi\right)^2}\right]}{1-\kappa \E\left[ \frac{1}{\left((\xi\Lambda)^{-1} +1\right)^2}\right]}.
\end{align}
{Again, note that this coincides with Equation \eqref{eq:gamma} that determines the parameter $\gamma$ in Definition \ref{def:Xi}, i.e., $\tau_*^2 = \gamma.$ }

\vp\noindent\underline{Proof of (iv):} 
For convenience, define 
$$F_n(\hat\betab,\betas,\Sigmab):=\frac{1}{p} \sum_{i=1}^{p} f\left(\sqrt{p}\hat\betab_{n,i},\sqrt{p}\betas_{i},\Sigmab_{ii}\right)\quad\text{and}\quad\alpha_*:=\E_\mu\left[f(X_{\kappa,\sigma^2}(\Lambda,B,H),B,\Lambda)\right].$$
Recall from \eqref{eq:w_n} the explicit expression for $\w_n$, repeated here for convenience.
\begin{align}
\w_n = -\left(\Sigmab^{-1}+\frac{u_n}{\tau_n}\Iden\right)^{-1}\left(\Sigmab^{-1/2}\betas-u_n\sqrt{\kappa}\hba\right).\nn
\end{align}
Also, recall that $\betab_n = \Sigmab^{-1/2}\w_n+\betas$. Thus,  (and using the fact that $\hba$ is distributed as $-\hba$),
\begin{align}
\betab_n &=\Sigmab^{-1/2}\left(\Sigmab^{-1}+\frac{u_n}{\tau_n}\Iden\right)^{-1}u_n\sqrt{\kappa}\hba + \left(\Iden-\Sigmab^{-1/2}\left(\Sigmab^{-1}+\frac{u_n}{\tau_n}\Iden\right)^{-1}\Sigmab^{-1/2}\right)\betas\nn\\
\Longrightarrow\betab_{n,i}&=\frac{{\Sigmab_{i,i}^{-1/2}}}{1+(\ksi_n\Sigmab_{i,i})^{-1}}\sqrt{\kappa}\tau_n\hba_i +\left(1-\frac{1}{1+\ksi_n\Sigmab_{i,i}}\right)\betas_i.
\end{align}
For $i\in[p]$, define 
\begin{align}
\vb_{n,i} = \frac{{\Sigmab_{i,i}^{-1/2}}}{1+(\ksi_*\Sigmab_{i,i})^{-1}}\sqrt{\kappa}\tau_*\hba_i +\left(1-\frac{1}{1+\ksi_*\Sigmab_{i,i}}\right)\betas_i
%-\left(\Sigmab^{-1}+\frac{u_*}{\tau_*}\Iden\right)^{-1}\left(\Sigmab^{-1/2}\betas-u_*\sqrt{\kappa}\hba\right).
\end{align}
In the above, for convenience, we have denoted $\xi_n:=u_n/\tau_n$ and recall that $\xi_*:=u_*/\tau_*$. 


The proof proceeds in two steps. In the first step, we use the fact that $\ksi_n\rP\ksi_*$ and $u_n\rP u_\star$ (see \eqref{eq:Dc}) to prove that for any $\eps\in(0,\ksi_*/2)$, there exists absolute constant $C>0$ such that wpa 1:
\begin{align}\label{eq:ivstep1}
|F_n(\sqrt{p}\betab_n,\sqrt{p}\betas,\Sigmab) - F_n(\sqrt{p}\vb_n,\sqrt{p}\betas,\Sigmab) | \leq C\eps.
\end{align}
In the second step, we use Lipschitzness of $f$ and Assumption \ref{ass:mu} to prove that 
\begin{align}
|F_n(\vb_n,\betas,\Sigmab)| \rP \alpha_*.\label{eq:ivstep2}
\end{align}
The desired follows by combining \eqref{eq:ivstep1} and \eqref{eq:ivstep2}. Thus, in what follows, we prove \eqref{eq:ivstep1} and \eqref{eq:ivstep2}.


\vspace{2pt}
\noindent\emph{Proof of \eqref{eq:ivstep1}.}~~Fix some $\eps\in(0,\ksi_*/2)$. From \eqref{eq:Dc}, we know that w.p.a. 1 $|\xi_n-\xi_*|\leq \eps$ and $|u_n-u_*|\leq \eps$. Thus, $\w_n$ is close to $\vb_n$. Specifically, in this event, for every $i\in[p]$, it holds that:
\begin{align}
\nn|\betab_{n,i} - \vb_{n,i}| &\leq {|\betas_i|}\left|\frac{1}{1+\xi_n\lai}-\frac{1}{1+\xi_*\lai}\right| + \sqrt{\kappa}\frac{|\hba_i|}{\sqrt{\Sigmab_{i,i}}}\left|\frac{\tau_n}{1+(\xi_n\lai)^{-1}}-\frac{\tau_*}{1+(\xi_*\lai)^{-1}}\right| \\
\nn&= {|\betas_i|}\left|\frac{1}{1+\xi_n\lai}-\frac{1}{1+\xi_*\lai}\right| + \sqrt{\kappa}\frac{|\hba_i|}{\sqrt{\Sigmab_{i,i}}}\left|\frac{u_n}{\xi_n+\lai^{-1}}-\frac{u_*}{\xi_*+\lai^{-1}}\right| \\
\nn& \leq
{|\betas_i|}\frac{|\lai||\xi_n-\xi_*|}{|1+\xi_n\lai||1+\xi_*\lai|} + \sqrt{\kappa}\frac{|\hba_i|}{\sqrt{\lai}}  \frac{u_*|\ksi_n-\ksi_*|}{(\ksi_n+\lai^{-1})(\ksi_*+\lai^{-1})} + 
\sqrt{\kappa}\frac{|\hba_i|}{\sqrt{\lai}}\frac{|u_n-u_*|}{\xi+\lai^{-1}}
\\
\nn& \leq
{|\betas_i|}{\Sigma_{\max}\eps} + \sqrt{\kappa}{|\hba_i|} u_*\Sigma_{\max}^{3/2} \eps+ 
\sqrt{\kappa}|\hba_i|\Sigma_{\max}^{1/2}\eps
\\
&\leq \eps \cdot \max\left\{\Sigma_{\max}^{3/2},\Sigma_{\max}^{1/2}\right\}\, \left(|\hba_i| + |\betas_i|\right).\label{eq:betav}
%\nn& \leq
%\frac{|\betas_i|}{\sqrt{\lai}} \frac{\eps}{(\xi-\eps)\xi} + \sqrt{\kappa}|\hba_i|\left(\frac{\eps(u_*+\eps)}{(\xi-\eps)\xi} + \frac{\eps}{\xi}\right)\\
%& \leq
%\frac{|\betas_i|}{\sqrt{\lai}} c_1(\eps) + |\hba_i|c_2(\eps).\label{eq:wivi}
\end{align}\som{$\eps$ missing}
In the second line above, we recalled that $u_n=\tau_n\xi_n$ and $u_*=\tau_*\xi_*$. In the third line, we used the triangle inequality. In the fourth line, we used that $\ksi_*>0$, $0<\lai\leq\Sigma_{\max}$ and $\ksi_n\geq \ksi_*-\eps \geq \ksi_*/2 >0$. 
%In the second line above, we used the triangle inequality. In the third inequality we used Assumption \ref{ass:inv} that $\lai>0, i\in[p]$, and in the last line we defined appropriate constants $c_1(\eps),c_2(\eps)>0$. Importantly,  ... 

Now, we will use this and Lipschitzness of $f$ to argue that there exists absolute constant $C>0$ such that wpa 1:
\begin{align}
|F_n(\sqrt{p}\betab_n,\sqrt{p}\betas,\Sigmab) - F_n(\sqrt{p}\vb_n,\sqrt{p}\betas,\Sigmab) | \leq C\eps.\nn
\end{align}
Denote, $\ab_i=(\sqrt{p}\betab_{n,i},\sqrt{p}\betas_i,\Sigmab_{i,i})$ and $\bb_i=(\sqrt{p}\vb_{n,i},\sqrt{p}\betas_i,\Sigmab_{i,i})$. Following the exact same argument as in \eqref{eq:LipC}, we have
\begin{align}
|F_n(\betab_n(\g,\h),\betas,\Sigmab) - F_n(\vb_n,\betas,\Sigmab)| &\leq \frac{L}{p}\sum_{i=1}^p (1+\max\{\|\ab_i\|_2^{k-1},\|\bb_i\|_2^{k-1}\})\|\ab_i-\bb_i\|_2\nn \\
&\leq \frac{L}{p}\sum_{i=1}^p \left(1+\max\{\|\ab_i\|_2^{k-1},\|\bb_i\|_2^{k-1}\}\right)\cdot|\sqrt{p}\betab_{n,i}(\g,\h) - \sqrt{p}\vb_{n,i}|\nn\\
&\leq L\left(1+\max\{\frac{1}{p}\sum_{i=1}^p \|\ab_i\|_2^{2k-2},\frac{1}{p}\sum_{i=1}^p\|\bb_i\|_2^{2k-2}\} \right)^{1/2}{\|\betab_{n}(\g,\h) - \vb_n\|_2}\nn.
\end{align}
From this and \eqref{eq:betav}, we find
\begin{align}
|F_n(\betab_n(\g,\h),\betas,\Sigmab) - F_n(\vb_n,\betas,\Sigmab)| &\leq 
L(1+\max\{\frac{1}{p}\sum_{i=1}^p \|\ab_i\|_2^{2k-2},\frac{1}{p}\sum_{i=1}^p\|\bb_i\|_2^{2k-2}\} )^{1/2} 
\\
&\qquad\qquad\eps\cdot\sqrt{2}\max\left\{\Sigma_{\max}^{3/2},\Sigma_{\max}^{1/2}\right\}\sqrt{\|\betas\|_2^2 + \|\hba\|_2^2}.\label{eq:epsS}
\end{align}
As in \eqref{eq:bdmom}, it can be shown that wpa 1: $\frac{1}{p}\sum_{i=1}^p \|\ab_i\|_2^{2k-2}<\infty$ and $\frac{1}{p}\sum_{i=1}^p \|\bb_i\|_2^{2k-2}<\infty$, as $p\rightarrow\infty$. Similarly, $\|\betas\|_2^2 = \frac{1}{p}\sum_{i=1}^p(\sqrt{p}\betas_i)^2<\infty$, as $p\rightarrow \infty$ {\color{red} by assumption on second moment convergence of $\sqrt{p}\betas$.} Finally, $\|\hba\|_2^2\leq 2$, wpa 1 as $p\rightarrow\infty$. Therefore, from \eqref{eq:epsS}, wpa 1, there exists constant $C>0$ such that 
\begin{align}
|F_n(\betab_n(\g,\h),\betas,\Sigmab) - F_n(\vb_n,\betas,\Sigmab)| \leq C \cdot\eps,
\end{align}
as desired.


%
%
%
% By $f\in\rm{PL}(k)$, we have that
%\begin{align}
%|f(\sqrt{p}\betab_{n,i},\sqrt{p}\betas_i,\Sigmab_{i,i}) - f(\sqrt{p}\vb_{n,i},\sqrt{p}\betas_i,\Sigmab_{i,i}) | \leq L ( 1 +  \|\ab_i\|_2^{k-1} +  \|\bb_i\|_2^{k-1} ) \sqrt{p} |\betab_{n,i}-\vb_{n,i}|. \nn
%\end{align}•
%Thus, using Cauchy-Swchartz,
%\begin{align}
%|F_n(\sqrt{p}\betab_n,\sqrt{p}\betas,\Sigmab) - F_n(\sqrt{p}\vb_n,\sqrt{p}\betas,\Sigmab) | &\leq \frac{L}{p} \sum_{i=1}^p ( 1 +  \|\ab_i\|_2^{k-1} +  \|\bb_i\|_2^{k-1} ) \sqrt{p} |\betab_{n,i}-\vb_{n,i}|. \\
%\end{align}•
%
% Using $f\in\rm{PL}(2)$, this implies the following
%\begin{align}
%|\frac{1}{p}\sum_{i=1}^{p}f(\sqrt{p}\w_{n,i})-\frac{1}{p}\sum_{i=1}^{p}f(\sqrt{p}\vb_{n,i})| &\leq \frac{1}{p}\sum_{i=1}^{p} L(1+\sqrt{p}|\w_{n,i}|+\sqrt{p}|\vb_{n,i}|)\sqrt{p}|\w_{n,i}-\vb_{n,i}| \nn
%\\
%&\leq \nn
%\frac{L}{p}\sum_{i=1}^{p}(1+\sqrt{p}|\w_{n,i}|+\sqrt{p}|\vb_{n,i}|) \left(|\nub_i| c_1 + |\h_i|c_2\right)\nn\\
%&\leq C\left(\frac{\tn{\nub}}{\sqrt{p}}+ \frac{\tn{\h}}{\sqrt{p}}\right)(1+B+\tn{\vb_{n}}) \label{eq:fwivi}
%%&\leq C(1+2B)\left(\frac{\tn{\nub}}{\sqrt{p}}+ \frac{\tn{\h}}{\sqrt{p}}\right) 
%\end{align}\som{$\eps$ missing}
%The inequality in the first line follows from $f\in\rm{PL}(2)$. The inequality in the second line uses \eqref{eq:wivi}, and we have also recalled the notation $\nub_i=\sqrt{p}\betas_i/\sqrt{\lai}$ in Assumption \ref{ass:mu} and $\h_i=\sqrt{p}\hba_i$. The inequality in the third line uses Cauchy-Swchartz, $\tn{\w_{n}}\leq B$ and $C:=C(L,c_1,c_2)$ is a universal constant.
%
%By Assumption \ref{ass:mu}, $\tn{\nub}/{\sqrt{p}}\rP\sqrt{\E[N^2]}<\infty$. Also, with probability approaching 1: $\h/{\sqrt{p}}<2$. Similarly, it can be argued that with probability approaching 1 $\tn{\w_n}$ is bounded by a large enough constant. Combined with \eqref{eq:fwivi}, we have shown that there exists universal constant $C'(\eps)>0$ such that with probability approaching 1:
%\begin{align}
%|\frac{1}{p}\sum_{i=1}^{p}f(\sqrt{p}\w_{n,i})-\frac{1}{p}\sum_{i=1}^{p}f(\sqrt{p}\vb_{n,i})| \leq C'(\eps) .
%\end{align}


\vspace{2pt}
\noindent\emph{Proof of \eqref{eq:ivstep2}.}~~Next, we will use Assumption \ref{ass:mu} to show that 
\begin{align}
|F_n(\vb_n,\betas,\Sigmab)| \rP \alpha_*.\label{eq:AO_conv}
\end{align}
Notice that $\vb_n$ is a function of $\betas,\Sigmab,\hba$. Concretely, define $g:\R^3\rightarrow\R$, such that
$$
g(x_1,x_2,x_3) := \frac{x_2^{-1/2}}{1+(\ksi x_2)^{-1}}\sqrt{\kappa}\tau_* x_3 + (1-(1+\ksi_*x_2)^{-1})x_1,
$$
and notice that
$$
\sqrt{p}\vb_{n,i} = g\left(\sqrt{p}\betas_i,\lai,\h_i\right) = \frac{{\Sigmab_{i,i}^{-1/2}}}{1+(\ksi_*\Sigmab_{i,i})^{-1}}\sqrt{\kappa}\tau_*\h_i +\left(1-\frac{1}{1+\ksi_*\Sigmab_{i,i}}\right)\sqrt{p}\betas_i.
$$
Thus,
$$
F_n(\vb_n,\betas,\Sigmab) = \frac{1}{p}\sum_{i=1}^p f\left(g\left(\sqrt{p}\betas_i,\lai,\h_i\right),\sqrt{p}\betas_i,\lai\right) =: \frac{1}{p}\sum_{i=1}^p h\left(\h_i,\sqrt{p}\betas_i,\lai\right),
$$
where we have defined $h:\R^3\rightarrow\R$:
\begin{align}
h(x_1,x_2,x_3) := f\left(g(x_2,x_3,x_1),x_2,x_3\right).\label{h func}
\end{align}
It will suffice to prove that $h\in\rm{PL}(k)$. Indeed, if that were the case, then Assumption \ref{ass:mu} would give
\begin{align}
\frac{1}{p}\sum_{i=1}^p h\left(\h_i,\sqrt{p}\betas_i,\lai\right) \rP \E_{\Nn(0,1)\otimes \mu}\left[h(H,B,\Lambda)\right] &=  \E_{\Nn(0,1)\otimes \mu}\left[f\left(g(B,\Lambda,H),B,\Lambda\right))\right]\\
& = \E[f\left(X_{\kappa,\sigma^2}(\Lambda,B,H),B,\Lambda\right)] = \alpha_*,
\end{align}
where the penultimate equality follows by recognizing that (cf. Eqn. \eqref{eq:X})
$$
g(B,\Lambda,H) = (1-(1+\ksi_*\Lambda)^{-1})B + \sqrt{\kappa}\frac{\tau_*\Lambda^{-1/2}}{1+(\ksi_*\Lambda)^{-1}}H =  X_{\kappa,\sigma^2}(\Lambda,B,H).
$$
In the remaining of the proof, \cts{we show that $h\in\rm{PL}(k)$} which is formalized below.

\begin{lemma} The function $h$ in \eqref{h func} is $\rm{PL}(k)$ as long as $f$ is $\rm{PL}(k)$.
\end{lemma}
\begin{proof} Since $f$ is $\rm{PL}(k)$, for some $L>0$, \eqref{PL func} holds. Fix $\ab=[\x,\y,\z]\in\R^{3p},\ab'=[\x',\y',\z']\in\R^{3p}$. Let $\bb=[\g,\y,\z]$ where $\g=g(\y,\z,\x)$. We have that
\begin{align}
|h(\ab)-h(\ab')|&\leq |f(\bb)-f(\bb')|\\
&\leq L\left(1+\|\bb\|_2^{k-1}+\|\bb'\|_2^{k-1}\right)\|\bb-\bb'\|_2\\
&\lesssim L\left(1+\|\ab\|_2^{k-1}+\|\ab'\|_2^{k-1}+\|\g\|_2^{k-1}+\|\g'\|_2^{k-1}\right)(\|\ab-\ab'\|_2+\|\g-\g'\|_2).
%~~~~~~~~&L\left(1+\|\g(\y,\z,\x)\|_2^{k-1}+\|g(\y',\z',\x')\|_2^{k-1}\right)\|g(\y,\z,\x)-g(\y',\z',\x')\|_2.
\end{align}
Next, we need to bound the $\g$ term in terms of $\ab$. This is accomplished as follows
\begin{align}
\|\g\|_2^{k-1}&\lesssim \frac{1}{p}\sum_{i=1}^p|g_i|^{k-1}\\
&\lesssim \frac{1}{p}\sum_{i=1}^p\left|\frac{y_i^{-1/2}}{1+(\ksi y_i)^{-1}}\sqrt{\kappa}\tau_* z_i + (1-(1+\ksi_*y_i)^{-1})x_i\right|\\
&\lesssim \frac{1}{p}\sum_{i=1}^p(|z_i|+|x_i|)^{k-1}\lesssim \|\x\|_2^{k-1}+\|\z\|_2^{k-1}\\
&\lesssim \|\ab\|^{k-1}.
\end{align}
Secondly and similarly, we have the following perturbation bound on the $\g-\g'$. \so{Suppose the triples $(x_i,y_i,z_i)$ takes values in a fixed bounded compact set $\Mc$. We will prove the following sequence of inequalities} 
\begin{align}
%|\|\g\|_2-\|\g'\|_2|&\leq 
\tn{\g-\g'}^2&\leq \sum_{i=1}^p|g(x_i,y_i,z_i)-g(x'_i,y'_i,z'_i)|^2\\
&\leq C^2_{\Mc} \sum_{i=1}^p(|x_i-x'_i|^2+|y_i-y'_i|^2+|z_i-z'_i|^2)\label{sec line eq}\\
&\leq C^2_{\Mc} (\tn{\x-\x'}^2+\tn{\y-\y'}^2+\tn{\z-\z'}^2)\\
&\lesssim \tn{\ab-\ab'}^2.
\end{align}
In what follows, we prove the second line \eqref{sec line eq} i.e.~the fact that for any triples $a=(x,y,z),a'=(x',y',z')$ (with $a,a'\in\R^3$), we have that 
\[
|g(a)-g(a')|\leq C_{\Mc}\tn{a-a'}.
\]
Set $C_\nabla=\sup_{a\in \Mc}\tn{\nabla g(a)}$. By definition of gradient, the inequality above holds with $C_\Mc=C_\nabla$. Thus, all that remains is proving that $C_\nabla$ is upper bounded by a constant. \so{However, this automatically holds because from the definition of $g(\dots)$ function, it is clear that $C_\nabla$ is defined everywhere and continuous thus it has a finite maximum over a compact set.}
% is bounded by a constant.
\end{proof}
%
%
%
%But from Assumption \ref{ass:mu} and $f\in\rm{PL(2)}$ it follows immediately that with probability approaching 1,
%\begin{align}
%|\frac{1}{p}\sum_{i=1}^{p}f(\sqrt{p}\vb_{n,i})-\E_\mu\left[X_{\kappa,\sigma^2}(\Lambda,N,H)  \right]| \leq \eps.
%\end{align}
%Combining the above displays and using triangle inequality completes the proof of the statement.
%


\section{Facts about Lipschitz functions}\label{SM useful fact}
For $k\geq 1$ we say a function $f:\R^m\rightarrow\R$ is pseudo-Lipschitz of order $k$ and denote it by $f\in \rm{PL}(k)$ if there exists a cosntant $L>0$ such that, for all $\x,\y\in\R^m$:
\begin{align}
|f(\x)-f(\y)|\leq L\left(1+\|\x\|_2^{k-1}+\|\y\|_2^{k-1}\right)\|\x-\y\|_2.\label{PL func}
\end{align}
In particular, when $f\in\rm{PL}(k)$, the following properties hold:
\begin{enumerate}
\item There exists a constant $L'$ such that for all $\x\in\R^n$: $|f(\x)|\leq L'(1+\|\x\|_2^k).$

\item $f$ is locally Lipschitz, that is for any $M>0$, there exists a constant $L_{M,m}<\infty$ such that for all $x,y\in[-M,M]^m$,
$
|f(\x)-f(\y)| \leq L_{M,m}\|\x-\y\|_2.
$
Further, $L_{M,m}\leq c(1+(M\sqrt{m})^{k-1})$ for some costant $c$.
\end{enumerate}


\begin{lemma}
Let $g:\R\rightarrow\R$ be a bounded, Lipschitz function. Consider the function $f:\R^3\rightarrow\R$ defined as follows:
$$
f(\x) = x_1(x_2-g(x_3))^2.
$$
Then, $f\in\rm{PL}(3).$
\end{lemma}

\begin{proof}
Let $h:\R^2\rightarrow\R$ defined as $h(\ub)=(\ub_1-g(\ub_2))^2$. The function $(\ub_1,\ub_2)\mapsto\ub_1-g(\ub_2)$ is clearly Lipschitz. Thus, $h\in\rm{PL}(2)$, i.e., 
\begin{align}
|h(\ub)-h(\vb)| \leq C(1+\|\ub\|_2+\|\vb\|_2)\|\ub-\vb\|_2\quad\text{and}\quad |h(\vb)|\leq C'(1+\|\vb\|_2^2).\label{eq:h_pl}
\end{align} 
Thus, letting $\x=(x_1,\ub)\in\R^3$ and $\y=(y_1,\vb)\in\R^3$, we have that 
\begin{align}
|f(\x)-f(\y)| &= |x_1h(\ub) - y_1h(\vb)| \leq |x_1||h(\ub)-h(\vb)| + |h(\vb)| |x_1-y_1|\nn\\
%&\leq |x_1|(1+\|\ub\|_2+\|\vb\|_2)\|\ub-\vb\|_2 + |h(\vb)| |x_1-y_1| \nn\\
&\leq C|x_1|(1+\|\ub\|_2+\|\vb\|_2)\|\ub-\vb\|_2 + C'(1+\|\vb\|_2^2)|x_1-y_1| \nn\\
&\leq C(|x_1|^2+(1+\|\ub\|_2+\|\vb\|_2)^2)\|\ub-\vb\|_2 + C'(1+\|\vb\|_2^2)|x_1-y_1| \nn\\
&\leq C(1+\|\x\|_2^2+\|\y\|_2^2)\|\ub-\vb\|_2 + C'(1+\|\x\|_2^2+\|\y\|_2^2)|x_1-y_1| \nn\\
&\leq C(1+\|\x\|_2^2+\|\y\|_2^2)\|\x-\y\|_2.
\end{align}
In the second line, we used \eqref{eq:h_pl}. In the third line, we used $2xy\leq x^2+y^2$. In the fourth line, we used Cauchy-Schwarz inequality. $C,C'>0$ are absolute constants that may change from line to line.

This completes the proof of the lemma.
\end{proof}



\begin{lemma}
Let functions $f,g:\R^3\rightarrow\R$ such that $f\in\rm{PL}(3)$ and 
$$
g(x_1,x_2,x_3) := \frac{x_2^{-1/2}}{1+(\ksi_* x_2)^{-1}}\sqrt{\kappa}\tau_* x_3 + (1-(1+\ksi_*x_2)^{-1})x_1.
$$
Here, $\ksi_*,\tau_*,\kappa$ are positive constants. Also, $x_2$ is bounded, say $x_2\in[m,M]\subset(0,\infty)$.
Further define 
\begin{align}
h(x_1,x_2,x_3) := f\left(g(x_2,x_3,x_1),x_2,x_3\right).
\end{align}
Then, it holds that $h\in\rm{PL}(k+?)$.
\end{lemma}

























